% part: intuitionistic-logic
% chapter: soundness-completeness
% section: soundness-axd

\documentclass[../../../include/open-logic-section]{subfiles}

\begin{document}

\olfileid{int}{sc}{sax}

\olsection{స్వీకృతాధారిత \usetoken{వ్యుత్పత్తుల}{derivation} నిర్దుష్టత}

\begin{editorial}
  నిర్దుష్టత నిరూపణ అన్ని స్వీకృతాలు అంతఃప్రజ్ఞావాద
  తర్కంలో చెల్లుబాటవుతాయనే వాస్తవంపై ఆధారపడుతుంది;
  దీన్ని ఇంకా నిరూపించాలి, ఉదాహరణకు అర్థవిజ్ఞాన
  అధ్యాయంలో.
\end{editorial}

\begin{thm}[నిర్దుష్టత]
  \ollabel{thm:soundness}
  $\Gamma \Proves !A$ అయితే, $ \Gamma \Entails !A$.
\end{thm}

\begin{proof}
  $\Gamma \Proves !A$ అయితే $\Gamma \Entails !A$
  అని నిరూపిస్తాం. $\Gamma$ నుంచి~$!A$ యొక్క
  \tetoken{వ్యుత్పత్తిలోని}{!!{derivation}} \tetoken{సూత్రాల}{!!{formula}s}
  సంఖ్య~$n$పై ఆగమన పద్ధతిలో నిరూపణ సాగుతుంది.
  $!A_1$, \dots, $!A_n = !A$ అనేది~$\Gamma$ నుంచి
  \tetoken{ఒక వ్యుత్పత్తి}{!!a{derivation}} అయితే $\Gamma
  \Entails !A_n$ అని చూపుతాం. $!A_1$, \dots, $!A_n$
  ఒక \tetoken{వ్యుత్పత్తి}{!!a{derivation}} అయితే, ఏ $k<n$కైనా
  $!A_1$, \dots, $!A_k$ కూడా వ్యుత్పత్తేనని గమనించండి.

  పొడవు~$0$ గల \tetoken{వ్యుత్పత్తులు}{!!{derivation}s}
  లేవు; కాబట్టి $n=0$కు వాదన శూన్యసత్యంగా నిలుస్తుంది.
  పొడవు~$<n$ గల అన్ని \tetoken{వ్యుత్పత్తులకు}{!!{derivation}s}
  వాదన నిలుస్తుందనుకుందాం. $!A_n$కు ఇచ్చిన
  కారణాన్ని బట్టి కేసులను వేరు చేస్తాం.

  \begin{enumerate}
  \item $!A_n$ ఒక స్వీకృతం. అన్ని స్వీకృతాలు
    చెల్లుబాటవుతాయి; కాబట్టి ఏ~$\Gamma$కైనా
    $\Gamma \Entails !A_n$.
  \item $!A_n \in \Gamma$. అప్పుడు ఏ
    $\mModel{M}$, $w$లకు అయినా $\mSat{M}{\Gamma}[w]$
    అయితే స్పష్టంగా $\mSat{M}{!A_n}[w]$;
    అంటే $\Gamma \Entails !A$.
    \sourcecorrection{OLTEINTSAX-001}{మూలంలో ఈ దశలో లోక-సత్య సంకేతానికి అదనపు $\{\Gamma\}$ వాదన జతైంది; నిర్వచిత రెండు వాదనల $\mSat{M}{!A_n}[w]$ రూపానికి సరిచేశాం.}
  \item $!A_i$, $!A_j \ident !A_i \lif !A_n$ నుంచి
    \MP{} ద్వారా $!A_n$ వస్తుంది. $!A_1$, \dots,
    $!A_i$ మరియు $!A_1$, \dots, $!A_j$
    అనేవి~$\Gamma$ నుంచి \tetoken{వ్యుత్పత్తులు}{!!{derivation}s}.
    అందువల్ల ఆగమన పరికల్పన ప్రకారం $\Gamma
    \Entails !A_i$, $\Gamma \Entails !A_i \lif !A_n$.

    $\mSat{M}{\Gamma}[w]$ అనుకుందాం.
    $\mSat{M}{\Gamma}[w]$, అలాగే
    $\Gamma \Entails !A_i \lif !A_n$ కలిపి
    $\mSat{M}{!A_i \lif
      !A_n}[w]$ ఇస్తాయి. నిర్వచనం ప్రకారం
    $Rww'$ అయ్యే ప్రతి $w'$ వద్ద
    $\mSat{M}{!A_i}[w']$ అయితే $\mSat{M}{!A_n}[w']$.
    $R$ స్వావర్తనం కాబట్టి, $Rww'$ అయ్యే
    $w'$లలో $w$ కూడా ఉంది; అంటే
    $\mSat{M}{!A_i}[w]$ అయితే $\mSat{M}{!A_n}[w]$.
    $\Gamma \Entails !A_i$ కాబట్టి
    $\mSat{M}{!A_i}[w]$. కనుక కోరినట్లుగా
    $\mSat{M}{!A_n}[w]$.
  \end{enumerate}
\end{proof}

\end{document}
